\chapter{Especificación de requisitos}

\section{Requisitos funcionales}

El sistema debe cubrir el proceso completo desde la descripción inicial del
problema hasta la generación de modelos ejecutables. Los requisitos funcionales
principales son:

\begin{table}[H]
\begin{center}
\begin{tabular}{|p{2.2cm}|p{10.2cm}|} \hline
\textbf{Código} & \textbf{Descripción} \\ \hline
RF-01 & Permitir iniciar una ejecución a partir de una descripción en lenguaje natural. \\ \hline
RF-02 & Identificar paradigmas científicos relevantes para el problema planteado. \\ \hline
RF-03 & Generar informes de investigación y guardarlos como artefactos persistentes. \\ \hline
RF-04 & Proponer formulaciones matemáticas para los paradigmas aprobados. \\ \hline
RF-05 & Solicitar una especificación del entorno de simulación antes de adaptar las formulaciones. \\ \hline
RF-06 & Generar especificaciones JSON implementables para el Builder. \\ \hline
RF-07 & Generar modelos Python autocontenidos y pruebas asociadas. \\ \hline
RF-08 & Ejecutar pruebas automáticas sobre los modelos generados. \\ \hline
RF-09 & Registrar modelos aprobados para su uso posterior en el laboratorio virtual. \\ \hline
RF-10 & Guardar en memoria solo artefactos aceptados tras revisión humana. \\ \hline
\end{tabular}
\caption{Requisitos funcionales principales.}
\end{center}
\end{table}

\section{Requisitos no funcionales}

Los requisitos no funcionales se centran en trazabilidad, robustez y
mantenibilidad.

\begin{itemize}
\item \textbf{Trazabilidad:} cada ejecución debe conservar estado, artefactos y
trazas de herramientas.
\item \textbf{Reanudación:} el sistema debe poder continuar una ejecución desde
el estado persistido.
\item \textbf{Modularidad:} cada agente debe tener una responsabilidad clara.
\item \textbf{Degradación controlada:} la generación principal debe seguir
funcionando aunque falten servicios opcionales de memoria.
\item \textbf{Validación:} los modelos generados deben acompañarse de pruebas
ejecutables.
\item \textbf{Separación de fases:} la salida de esta fase debe integrarse con
la fase de simulación sin acoplamiento rígido.
\end{itemize}

\section{Información almacenada}

El sistema almacena información en varias capas:

Esta separación sigue la lógica de los sistemas de recuperación aumentada y de
memoria externa: los artefactos completos se conservan como fuente auditable,
los hechos se indexan para recuperación densa y léxica, y las relaciones se
mantienen en un grafo de conocimiento
\cite{lewis2020rag,karpukhin2020dpr,robertson2009bm25,hogan2021knowledgegraphs,zhong2023memorybank,packer2023memgpt}.

\begin{table}[H]
\begin{center}
\begin{tabular}{|p{3cm}|p{9.5cm}|} \hline
\textbf{Sistema} & \textbf{Información almacenada} \\ \hline
MinIO & Informes, formulaciones, especificaciones, código generado, pruebas, estado y trazas. \\ \hline
PostgreSQL & Ejecuciones, metadatos de artefactos, modelos registrados y ciclo de vida de memorias. \\ \hline
Neo4j & Paradigmas, papers, variables, postulados, formulaciones, modelos y relaciones. \\ \hline
Qdrant & Hechos de memoria indexados mediante embeddings densos y BM25. \\ \hline
\end{tabular}
\caption{Información persistida por el sistema.}
\end{center}
\end{table}

\section{Interfaces con otros sistemas}

La fase desarrollada se comunica con:

\begin{itemize}
\item proveedores de modelos de lenguaje;
\item servicios de búsqueda web o académica;
\item infraestructura compartida del repositorio;
\item herramientas de ejecución de pruebas;
\item la fase de simulación, mediante el contrato de modelo generado.
\end{itemize}

La interfaz más relevante con la segunda fase es el contrato
\texttt{DecisionModel}. Este contrato evita que los modelos generados dependan
de clases internas del simulador.
